\MathBookSourcePage{114}
\setcounter{statement}{2}
\begin{Lemma}\label{lem:appendix-moment-interpolant-local}
设 $p\in\mathbb{R}$ 且 $p>1$.
对每个 $v\in W^{2,p}(\kappa)$, 恰有一个多项式
$\widetilde I_\kappa v\in P_k$ 满足
\begin{equation}\label{eq:appendix-moment-interpolant-local}
\left\{
\begin{aligned}
&\widetilde I_\kappa v(a_i)=v(a_i),
&&1\leq i\leq3,\\
&\text{若 }k\geq2,
&&\int_{\kappa'}(\widetilde I_\kappa v-v)f\,\mathrm{d}s=0
\quad
\forall f\in P_{k-2}(\kappa'),
\quad
\text{对 $\kappa$ 的每条边 $\kappa'$},\\
&\text{若 }k\geq3,
&&\int_\kappa(\widetilde I_\kappa v-v)f\,\mathrm{d}x=0
\quad
\forall f\in P_{k-3}(\kappa).
\end{aligned}
\right.
\end{equation}
\end{Lemma}
\begin{Proof}
首先注意, \eqref{eq:appendix-moment-interpolant-local} 包含
$\frac12(k+1)(k+2)$ 个方程, 这恰好是 $P_k$ 的维数.
因此, \eqref{eq:appendix-moment-interpolant-local} 是一个方形线性方程组,
只需证明其解唯一. 为此假设 $p\in P_k$ 满足
\begin{align*}
&\text{(i)}\quad p(a_i)=0,
&&1\leq i\leq3,\\
&\text{(ii)}\quad
\int_{\kappa'}p(s)f(s)\,\mathrm{d}s=0,
&&\forall f\in P_{k-2}(\kappa'),
\quad\text{对 $\kappa$ 的每条边 $\kappa'$},\\
&\text{(iii)}\quad
\int_\kappa p(x)f(x)\,\mathrm{d}x=0,
&&\forall f\in P_{k-3}(\kappa).
\end{align*}
先设 $k\geq3$. 注意 $p$ 在每条 $\kappa'$ 上的限制是单变量 $s$
的 $k$ 次多项式. 由 (i) 和 (ii),
\[
0
=
\int_{\kappa'}p(s)\frac{\mathrm{d}^2p(s)}{\mathrm{d}s^2}\,\mathrm{d}s
=
-\int_{\kappa'}
\left(\frac{\mathrm{d}p(s)}{\mathrm{d}s}\right)^2\,\mathrm{d}s.
\]
因此 $p=0$ 于 $\partial\kappa$, 且可用重心坐标表示为
\[
p=\lambda_1\lambda_2\lambda_3q,
\qquad
q\in P_{k-3}.
\]
于是 (iii) 蕴含
\[
\int_\kappa\lambda_1\lambda_2\lambda_3q^2(x)\,\mathrm{d}x=0.
\]
被积函数非负, 故 $q=0$, 从而 $p=0$ 于 $\kappa$.
$k\leq2$ 的情形是平凡的.
\end{Proof}

按常规容易验证
$\widetilde I_\kappa\in\mathcal{L}(W^{2,p}(\kappa);P_k)$,
$\widetilde I_\kappa v$ 在 $\kappa$ 的每条边 $\kappa'$ 上的值只依赖于
$v$ 在 $\kappa'$ 上的值, 并且
\[
\widetilde I_\kappa p=p
\qquad
\forall p\in P_k.
\]
此外, 算子 $\widetilde I_\kappa$ 满足基本关系
\eqref{eq:appendix-interpolant-commuting-definition}:
\[
\widehat{\widetilde I_\kappa v}
=
\widetilde I_{\hat\kappa}\hat v.
\]
汇集这些性质, 并应用 Corollary
\ref{cor:appendix-abstract-interpolation-physical} 和 Remark
\ref{rem:appendix-local-interpolation-estimate}, 得到:

\MathBookSourcePage{115}
\setcounter{statement}{3}
\begin{Lemma}\label{lem:appendix-moment-interpolant-global}
设 $\mathcal{T}_h$ 是 $\bar\Omega$ 的一个正则三角剖分.
对实数 $p>1$, 在每个单元 $\kappa$ 上由
\[
\widetilde I_hv|_\kappa
=
\widetilde I_\kappa v
\qquad
\forall\kappa\in\mathcal{T}_h
\]
定义的插值算子
\[
\widetilde I_h
\in
\mathcal{L}(W^{2,p}(\Omega);\Theta_h)
\cap
\mathcal{L}(W^{2,p}(\Omega)\cap W_0^{1,p}(\Omega);\Phi_h)
\]
对所有满足 $0\leq m\leq s+1$, $1\leq s\leq k$ 的整数 $m$
和实数 $s$ 有误差估计
\begin{equation}\label{eq:appendix-moment-interpolation-error}
|v-\widetilde I_hv|_{m,p,\Omega}
\leq
Ch^{s+1-m}|v|_{s+1,p,\Omega}
\qquad
\forall v\in W^{s+1,p}(\Omega),
\end{equation}
其中常数 $C>0$ 与 $h$ 和 $v$ 无关.
\end{Lemma}

除了这两个插值算子以外, 投影算子将在后续理论中起基本作用.
更确切地, 对实数 $p\geq1$, 令
\[
\dot P_h\in\mathcal{L}(W_0^{1,p}(\Omega);\Phi_h),
\qquad
P_h\in\mathcal{L}(W^{1,p}(\Omega);\Theta_h)
\]
分别由
\begin{equation}\label{eq:appendix-dirichlet-ritz-projection}
(\operatorname{grad}(\dot P_hv-v),
\operatorname{grad}\phi_h)=0
\quad
\forall\phi_h\in\Phi_h,
\quad
\forall v\in W_0^{1,p}(\Omega)
\end{equation}
和
\begin{equation}\label{eq:appendix-neumann-ritz-projection}
\left\{
\begin{aligned}
(\operatorname{grad}(P_hv-v),
\operatorname{grad}\theta_h)&=0
&&\forall\theta_h\in\Theta_h,
\quad\forall v\in W^{1,p}(\Omega),\\
(P_hv-v,1)&=0
\end{aligned}
\right.
\end{equation}
定义.

注意, $\dot P_hv$ 也可以解释为齐次 Dirichlet 问题
\[
-\Delta v=f\quad\text{in }\Omega,
\qquad
v=0\quad\text{on }\Gamma
\]
的有限元解, 而 $P_hv$ 可以解释为非齐次 Neumann 问题
\[
-\Delta v=f\quad\text{in }\Omega,
\qquad
\frac{\partial v}{\partial n}=g\quad\text{on }\Gamma
\]
的有限元解, 其中相容条件为
\[
\int_\Omega f\,\mathrm{d}x
=
\langle g,1\rangle_\Gamma,
\]
其唯一性由 \eqref{eq:appendix-neumann-ritz-projection} 的第二个条件得到.
由于 $\dot P_h$ 和 $P_h$ 是关于 $H^1(\Omega)$ 半范数的投影,
当 $v\in H^1(\Omega)$ 时, 容易证明 $L^2$ 和 $H^1$ 范数中的误差估计;
但要建立最优 $L^p$ 和 $W^{1,p}$ 估计则困难得多.
下面的 Theorem 是多位数学家多年工作的成果, 例如参见
Douglas, Dupont \& Wahlbin [25], Nitsche [62],
Rannacher \& Scott [66] 以及 Scott [73].
虽然它绝非标准结果, 但其证明远超出本书范围, 因而从略.

\setcounter{statement}{1}
\begin{Theorem}\label{thm:appendix-ritz-projection-lp-error}
假设 $\Omega$ 是凸多边形. 设 $\mathcal{T}_h$ 是 $\bar\Omega$
的一个一致正则三角剖分, 实数 $s$ 和 $p$ 满足
$0\leq s\leq k$ 以及 $1\leq p\leq\infty$.
当 $k\geq2$, 或者 $k=1$ 且 $2\leq p<\infty$ 时,
存在与 $h$ 无关的常数 $C>0$, 使投影 $P_h$(相应地, $\dot P_h$)
满足误差估计
\begin{equation}\label{eq:appendix-ritz-projection-lp-error}
\begin{aligned}
&\|v-P_hv\|_{0,p,\Omega}
+h|v-P_hv|_{1,p,\Omega}
\leq
Ch^{s+1}\|v\|_{s+1,p,\Omega},\\
&\forall v\in W^{s+1,p}(\Omega)
\quad
(\text{相应地, }
\forall v\in W^{s+1,p}(\Omega)\cap W_0^{1,p}(\Omega)).
\end{aligned}
\end{equation}
当 $k=1$ 且 $p\in[1,2)$ 时, 估计
\eqref{eq:appendix-ritz-projection-lp-error} 的右端还要乘以因子
\[
|\ln h|^{2/p-1}.
\]
当 $k=1$ 且 $p=\infty$ 时,
\eqref{eq:appendix-ritz-projection-lp-error} 中的 $L^\infty$ 估计变为
\begin{equation}\label{eq:appendix-ritz-projection-linfty-error}
\|v-P_hv\|_{0,\infty,\Omega}
\leq
C|\ln h|h^{s+1}\|v\|_{s+1,\infty,\Omega},
\end{equation}
而 $W^{1,\infty}$ 估计保持不变.
\end{Theorem}

\MathBookSourcePage{116}
Theorem \ref{thm:appendix-ritz-projection-lp-error} 的证明依赖于如下事实:
对某个 $\varepsilon>0$ 和所有 $p\in(1,2+\varepsilon)$,
Laplacian 算子是从
$W_0^{1,p}(\Omega)\cap W^{2,p}(\Omega)$ 到 $L^p(\Omega)$ 的同构.
根据 Remark \ref{rem:laplace-isomorphism-convex-domain},
这个同构在光滑区域或凸多边形上成立.
这解释了 Theorem \ref{thm:appendix-ritz-projection-lp-error}
中的凸性假设. 还应指出, 对数因子的出现是 $L^\infty$ 估计中的一个
众所周知的现象. 最后, 如上所述, 当 $p=2$ 时,
\eqref{eq:appendix-ritz-projection-lp-error} 可对所有 $k\geq1$ 直接证明;
若 $s$ 是整数, 则得到 $P_h$(相应地, $\dot P_h$) 的熟知估计
\begin{equation}\label{eq:appendix-ritz-projection-h1-error}
\begin{aligned}
&\|v-P_hv\|_{0,\Omega}
+h|v-P_hv|_{1,\Omega}
\leq
Ch^{m+1}|v|_{m+1,\Omega},\\
&\forall v\in H^{m+1}(\Omega)
\quad
(\text{相应地, }
\forall v\in H^{m+1}(\Omega)\cap H_0^1(\Omega)).
\end{aligned}
\end{equation}

\setcounter{statement}{1}
\begin{Remark}\label{rem:appendix-discontinuous-function-projections}
当函数 $v$ 不连续时, 它的插值 $I_hv$ 和 $\widetilde I_hv$ 都没有定义.
若 $\Omega$ 是凸的且三角剖分一致正则, 可方便地用投影 $P_hv$
替代它们(若 $v$ 在 $\Gamma$ 上为零, 则用 $\dot P_hv$).
当 $\Omega$ 非凸或没有一致正则三角剖分时,
可以使用由局部正则化得到的其他插值算子.
与投影相比, 它们还有局部而非全局定义的优点.
关于这种技术的一些内容见 Section
\ref{sec:appendix-interpolation-discontinuous-functions}.
\end{Remark}

后文还会使用局部 $L^2$ 投影. 更确切地, 对
$v\in L^2(\Omega)$ 和每个整数 $k\geq0$, 定义
\begin{equation}\label{eq:appendix-local-l2-projection}
\rho_hv|_\kappa\in P_k,
\qquad
\int_\kappa(\rho_hv-v)f\,\mathrm{d}x=0
\quad
\forall f\in P_k,
\quad
\forall\kappa\in\mathcal{T}_h.
\end{equation}
显然, 算子 $\rho_h$ 满足 Corollary
\ref{cor:appendix-abstract-interpolation-physical} 的假设,
从而得到下面在所有维数 $N$ 中均成立的结果.

\setcounter{statement}{4}
\begin{Lemma}\label{lem:appendix-local-l2-projection-error}
设 $v\in H^s(\Omega)$, 其中实数 $s\in[0,k+1]$.
则 $L^2$ 投影 $\rho_h$ 满足误差估计
\begin{equation}\label{eq:appendix-local-l2-projection-error}
\|v-\rho_hv\|_{0,\Omega}
\leq
Ch^s|v|_{s,\Omega},
\end{equation}
其中常数 $C>0$ 与 $h$ 和 $v$ 无关.
\end{Lemma}
注意, 这个 Lemma 不要求三角剖分具有正则性.

\MathBookSourcePage{117}
本节最后给出 $\Theta_h$ 中函数在任意维数 $N$ 下满足的若干逆不等式.

\setcounter{statement}{5}
\begin{Lemma}\label{lem:appendix-local-inverse-inequality}
设 $r$ 和 $p$ 为满足 $1\leq r,p\leq\infty$ 的实数.
若 $1/r-1/p\geq1/N$, 假设 $\mathcal{T}_h$ 正则;
否则假设 $\mathcal{T}_h$ 一致正则.
则存在与 $h$ 和 $\kappa$ 无关的常数 $C>0$, 使得
\begin{equation}\label{eq:appendix-local-inverse-inequality}
|v|_{1,r,\kappa}
\leq
Ch^{N(1/r-1/p)-1}\|v\|_{0,p,\kappa}
\qquad
\forall\kappa\in\mathcal{T}_h,
\quad
\forall v\in\Theta_h.
\end{equation}
\end{Lemma}
\begin{Proof}
设 $v\in\Theta_h$. 由
\eqref{eq:appendix-change-inverse},
\eqref{eq:appendix-simplex-matrix-bounds} 和
\eqref{eq:appendix-simplex-determinant-measure},
\[
|v|_{1,r,\kappa}
\leq
C_1\rho_\kappa^{-1}
(\operatorname{meas}(\kappa))^{1/r}
|\hat v|_{1,r,\hat\kappa}.
\]
由于 $\hat v$ 属于 $\hat\kappa$ 上的有限维空间 $P_k$, 有
\[
|\hat v|_{1,r,\hat\kappa}
\leq
C_2\|\hat v\|_{0,p,\hat\kappa},
\]
其中常数 $C_2$ 仅依赖于 $r$, $p$, $k$, $N$ 和 $\hat\kappa$ 的几何.
再应用 \eqref{eq:appendix-change-forward} 和
\eqref{eq:appendix-simplex-determinant-measure}, 得
\begin{equation}\label{eq:appendix-local-inverse-intermediate}
|v|_{1,r,\kappa}
\leq
C_3\rho_\kappa^{-1}
(\operatorname{meas}(\kappa))^{1/r-1/p}
\|v\|_{0,p,\kappa}.
\end{equation}
根据 \eqref{eq:appendix-simplex-determinant-bounds},
若 $N(1/r-1/p)\geq1$,
\eqref{eq:appendix-local-inverse-intermediate} 的右端含 $h_\kappa$
的正幂, 因而 $\mathcal{T}_h$ 的正则性足以推出
\eqref{eq:appendix-local-inverse-inequality}.
否则, 要得到相同结论就必须要求 $\mathcal{T}_h$ 一致正则.
\end{Proof}

\setcounter{statement}{2}
\begin{Corollary}\label{cor:appendix-global-inverse-gradient}
设 $r$ 和 $p$ 如 Lemma \ref{lem:appendix-local-inverse-inequality}
中所述, 并假设 $\mathcal{T}_h$ 是 $\Omega$ 的一个一致正则三角剖分.
则存在与 $h$ 无关的常数 $C>0$, 使得
\begin{equation}\label{eq:appendix-global-inverse-gradient}
|v|_{1,r,\Omega}
\leq
Ch^{-1+\min(0,N/r-N/p)}
\|v\|_{0,p,\Omega}
\qquad
\forall v\in\Theta_h.
\end{equation}
\end{Corollary}
\begin{Proof}
若 $r\geq p$, 则由
\eqref{eq:appendix-local-inverse-inequality} 和 Jensen 不等式
\begin{equation}\label{eq:appendix-discrete-jensen}
\left(\sum_{i=1}^I|a_i|^r\right)^{1/r}
\leq
\left(\sum_{i=1}^I|a_i|^p\right)^{1/p},
\end{equation}
得到 \eqref{eq:appendix-global-inverse-gradient};
该不等式对每个有限和都成立.

若 $r<p$, 使用 H\"older 不等式的离散形式
\begin{equation}\label{eq:appendix-discrete-holder}
\left(\sum_{i=1}^I|a_i|^r|b_i|^{1-r/p}\right)^{1/r}
\leq
\left(\sum_{i=1}^I|a_i|^p\right)^{1/p}
\left(\sum_{i=1}^I|b_i|\right)^{1/r-1/p},
\end{equation}
其中 $|a_i|=\|v\|_{0,p,\kappa}$,
$b_i=\operatorname{meas}(\kappa)$, 求和遍历 $\mathcal{T}_h$ 的所有
$\kappa$. 因此由 $\mathcal{T}_h$ 的一致正则性和
\eqref{eq:appendix-local-inverse-intermediate},
\[
|v|_{1,r,\Omega}
\leq
C_3\frac{\sigma}{\tau h}
(\operatorname{meas}(\Omega))^{1/r-1/p}
\|v\|_{0,p,\Omega}.
\]
\end{Proof}

\setcounter{statement}{6}
\begin{Lemma}\label{lem:appendix-global-inverse-same-order}
设 $r$ 和 $p$ 如 Lemma \ref{lem:appendix-local-inverse-inequality}
中所述, $m$ 为非负整数. 若 $r\leq p$, 或者当 $r>p$ 时
$\mathcal{T}_h$ 一致正则, 则存在与 $h$ 无关的常数 $C>0$, 使得
\begin{equation}\label{eq:appendix-global-inverse-same-order}
|v|_{m,r,\Omega}
\leq
Ch^{\min(0,N/r-N/p)}|v|_{m,p,\Omega}
\qquad
\forall v\in\Theta_h.
\end{equation}
\end{Lemma}
\begin{Proof}
当 $r\leq p$ 时, \eqref{eq:appendix-global-inverse-same-order} 显然成立.
当 $r>p$ 时, 证明与
\eqref{eq:appendix-global-inverse-gradient} 的证明十分相似, 留给读者.
\end{Proof}

\MathBookSourcePage{118}
\section{四边形有限元}\label{sec:appendix-quadrilateral-finite-elements}

本节专门讨论平面多边形区域. 除非另有说明, 记号沿用上一节.
此外, 由于三角形有限元的若干性质对四边形有限元仍然成立,
我们将重点考察四边形所特有的性质.

\setcounter{statement}{2}
\begin{Definition}\label{def:appendix-polynomial-space-qk}
对每个整数 $k\geq0$, 记 $Q_k$ 为参考空间
$(\hat x_1,\hat x_2)$ 中形如
\[
q(\hat x)=\sum c_{ij}\hat x_1^i\hat x_2^j
\]
的所有多项式构成的空间, 其中求和遍历满足 $0\leq i,j\leq k$
的所有整数 $i,j$.
\end{Definition}

当然 $Q_0=P_0$, 但对所有 $k\geq1$ 都有严格包含
\[
P_k\subset Q_k.
\]
